\documentclass[10pt,a4paper]{amsart}
\usepackage{amsmath,amssymb,amsthm,graphicx}
% Retained preamble for the Stacks Project chapter divisors.tex (Divisors).
% The chapter file pinned-original-source.tex carries only \input{preamble} on its first line, so no
% preamble text lives inside the chapter itself.  The Stacks Project uses one single shared preamble for
% every chapter; the snapshot stages it verbatim in sources/stacks-chow/preamble.tex.  The block below is
% copied mechanically, line for line, from that file's lines 14-18 (the Xy-pic setup) and lines 154-217
% (the native theorem declarations together with the mathematical macros).  Nothing was added, removed or
% redefined.  The wrapper class and the amsmath/amssymb/amsthm/graphicx base packages are supplied by the
% assembler, so the class/\externaldocument part of the shared preamble is not repeated here.
% For commutative diagrams we use Xy-pic
\usepackage[all]{xy}

% We use 2cell for 2-commutative diagrams.
\xyoption{2cell}
% Theorem environments.
%
\theoremstyle{plain}
\newtheorem{theorem}[subsection]{Theorem}
\newtheorem{proposition}[subsection]{Proposition}
\newtheorem{lemma}[subsection]{Lemma}

\theoremstyle{definition}
\newtheorem{definition}[subsection]{Definition}
\newtheorem{example}[subsection]{Example}
\newtheorem{exercise}[subsection]{Exercise}
\newtheorem{situation}[subsection]{Situation}

\theoremstyle{remark}
\newtheorem{remark}[subsection]{Remark}
\newtheorem{remarks}[subsection]{Remarks}

\numberwithin{equation}{subsection}

% Macros
%
\def\lim{\mathop{\mathrm{lim}}\nolimits}
\def\colim{\mathop{\mathrm{colim}}\nolimits}
\def\Spec{\mathop{\mathrm{Spec}}}
\def\Hom{\mathop{\mathrm{Hom}}\nolimits}
\def\Ext{\mathop{\mathrm{Ext}}\nolimits}
\def\SheafHom{\mathop{\mathcal{H}\!\mathit{om}}\nolimits}
\def\SheafExt{\mathop{\mathcal{E}\!\mathit{xt}}\nolimits}
\def\Sch{\mathit{Sch}}
\def\Mor{\mathop{\mathrm{Mor}}\nolimits}
\def\Ob{\mathop{\mathrm{Ob}}\nolimits}
\def\Sh{\mathop{\mathit{Sh}}\nolimits}
\def\NL{\mathop{N\!L}\nolimits}
\def\CH{\mathop{\mathrm{CH}}\nolimits}
\def\proetale{{pro\text{-}\acute{e}tale}}
\def\etale{{\acute{e}tale}}
\def\QCoh{\mathit{QCoh}}
\def\Ker{\mathop{\mathrm{Ker}}}
\def\Im{\mathop{\mathrm{Im}}}
\def\Coker{\mathop{\mathrm{Coker}}}
\def\Coim{\mathop{\mathrm{Coim}}}

% Boxtimes
%
\DeclareMathSymbol{\boxtimes}{\mathbin}{AMSa}{"02}

%
% Macros for moduli stacks/spaces
%
\def\QCohstack{\mathcal{QC}\!\mathit{oh}}
\def\Cohstack{\mathcal{C}\!\mathit{oh}}
\def\Spacesstack{\mathcal{S}\!\mathit{paces}}
\def\Quotfunctor{\mathrm{Quot}}
\def\Hilbfunctor{\mathrm{Hilb}}
\def\Curvesstack{\mathcal{C}\!\mathit{urves}}
\def\Polarizedstack{\mathcal{P}\!\mathit{olarized}}
\def\Complexesstack{\mathcal{C}\!\mathit{omplexes}}
% \Pic is the operator that assigns to X its picard group, usage \Pic(X)
% \Picardstack_{X/B} denotes the Picard stack of X over B
% \Picardfunctor_{X/B} denotes the Picard functor of X over B
\def\Pic{\mathop{\mathrm{Pic}}\nolimits}
\def\Picardstack{\mathcal{P}\!\mathit{ic}}
\def\Picardfunctor{\mathrm{Pic}}
\def\Deformationcategory{\mathcal{D}\!\mathit{ef}}

% --- appended 2026-10-12 (Chief): the project own definitions of the reference/slogan/history
% display environments, copied mechanically from the shared Stacks preamble lines 8-12
% (\usepackage{verbatim} plus the three \newenvironment lines, which the chapter preamble keeps
% verbatim from the project preamble). Nothing added, removed or redefined: these environments
% carry no mathematical content and the project itself typesets them as comments. Without these
% three lines a counted statement whose own text uses them fails with Environment slogan undefined.
\usepackage{verbatim}
\newenvironment{reference}{\comment}{\endcomment}
\newenvironment{slogan}{\comment}{\endcomment}
\newenvironment{history}{\comment}{\endcomment}
\usepackage{cleveref}
\crefname{originallemma}{original lemma}{original lemmas}
\Crefname{originallemma}{Original lemma}{Original lemmas}
\makeatletter
\newlabel{algebra-lemma-characterize-dvr}{{00PD}{source}{}{sourcepointer-1}{}}
\newlabel{algebra-lemma-characterize-dvr@cref}{{[originallemma][0][]00PD}{[1][0][]source}{}{}{}}
\makeatother
\crefname{originallemma}{original lemma}{original lemmas}
\Crefname{originallemma}{Original lemma}{Original lemmas}
\makeatletter
\newlabel{algebra-lemma-criterion-normal}{{031S}{source}{}{sourcepointer-2}{}}
\newlabel{algebra-lemma-criterion-normal@cref}{{[originallemma][0][]031S}{[1][0][]source}{}{}{}}
\makeatother
\crefname{originallemma}{original lemma}{original lemmas}
\Crefname{originallemma}{Original lemma}{Original lemmas}
\makeatletter
\newlabel{algebra-lemma-flat-intersect-ideals}{{0BBY}{source}{}{sourcepointer-3}{}}
\newlabel{algebra-lemma-flat-intersect-ideals@cref}{{[originallemma][0][]0BBY}{[1][0][]source}{}{}{}}
\makeatother
\crefname{originallemma}{original lemma}{original lemmas}
\Crefname{originallemma}{Original lemma}{Original lemmas}
\makeatletter
\newlabel{algebra-lemma-minimal-over-1}{{00KV}{source}{}{sourcepointer-4}{}}
\newlabel{algebra-lemma-minimal-over-1@cref}{{[originallemma][0][]00KV}{[1][0][]source}{}{}{}}
\makeatother
\crefname{originallemma}{original lemma}{original lemmas}
\Crefname{originallemma}{Original lemma}{Original lemmas}
\makeatletter
\newlabel{algebra-lemma-normal-domain-intersection-localizations-height-1}{{031T}{source}{}{sourcepointer-5}{}}
\newlabel{algebra-lemma-normal-domain-intersection-localizations-height-1@cref}{{[originallemma][0][]031T}{[1][0][]source}{}{}{}}
\makeatother
\crefname{originallemma}{original lemma}{original lemmas}
\Crefname{originallemma}{Original lemma}{Original lemmas}
\makeatletter
\newlabel{algebra-lemma-product-ideals-principal}{{09ME}{source}{}{sourcepointer-6}{}}
\newlabel{algebra-lemma-product-ideals-principal@cref}{{[originallemma][0][]09ME}{[1][0][]source}{}{}{}}
\makeatother
\crefname{originallemma}{original lemma}{original lemmas}
\Crefname{originallemma}{Original lemma}{Original lemmas}
\makeatletter
\newlabel{algebra-lemma-symbolic-power-flat-extension}{{0BC0}{source}{}{sourcepointer-7}{}}
\newlabel{algebra-lemma-symbolic-power-flat-extension@cref}{{[originallemma][0][]0BC0}{[1][0][]source}{}{}{}}
\makeatother
\crefname{originallemma}{original lemma}{original lemmas}
\Crefname{originallemma}{Original lemma}{Original lemmas}
\makeatletter
\newlabel{lemma-flat-over-integral-integral-fibre}{{0BCM}{source}{}{sourcepointer-8}{}}
\newlabel{lemma-flat-over-integral-integral-fibre@cref}{{[originallemma][0][]0BCM}{[1][0][]source}{}{}{}}
\makeatother
\crefname{originallemma}{original lemma}{original lemmas}
\Crefname{originallemma}{Original lemma}{Original lemmas}
\makeatletter
\newlabel{lemma-pullback-meromorphic-sections-defined}{{02OU}{source}{}{sourcepointer-9}{}}
\newlabel{lemma-pullback-meromorphic-sections-defined@cref}{{[originallemma][0][]02OU}{[1][0][]source}{}{}{}}
\makeatother
\crefname{originallemma}{original lemma}{original lemmas}
\Crefname{originallemma}{Original lemma}{Original lemmas}
\makeatletter
\newlabel{morphisms-lemma-fppf-open}{{01UA}{source}{}{sourcepointer-10}{}}
\newlabel{morphisms-lemma-fppf-open@cref}{{[originallemma][0][]01UA}{[1][0][]source}{}{}{}}
\makeatother
\crefname{originallemma}{original lemma}{original lemmas}
\Crefname{originallemma}{Original lemma}{Original lemmas}
\makeatletter
\newlabel{morphisms-lemma-fpqc-quotient-topology}{{02JY}{source}{}{sourcepointer-11}{}}
\newlabel{morphisms-lemma-fpqc-quotient-topology@cref}{{[originallemma][0][]02JY}{[1][0][]source}{}{}{}}
\makeatother
\title{An invertible module with trivial generic fibre over a flat family with integral fibres on a normal base is pulled back from the base}
\author{The Stacks Project Authors}
\date{stacks-project:divisors.tex:local-snapshot-20261012; GNU FDL 1.2 or later}
\begin{document}
\maketitle
\noindent\textit{Editorial note.} All delivered mathematical text is the author's own, copied from the pinned Stacks Project chapter file pinned-original-source.tex; no statement or proof step was rewritten, reordered or removed.\par\medskip
\noindent\textit{Editorial note.} Following the multi-claim convention for this source, source\_url is the source-level URL https://stacks.math.columbia.edu/ and the source label lemma-in-image-pullback is carried in claim\_key and here rather than in a per-claim tag URL.\par\medskip
\noindent\textit{Editorial note.} The preamble is the chapter's shared Stacks Project preamble, retained as sources/stacks-divisors/retained-preamble.tex: a line-for-line excerpt of the shared Stacks preamble staged with this snapshot. No macro was invented or redefined, and this card adds no preamble text of its own.\par\medskip
\noindent\textit{Editorial note.} Statement span 6967-6981 and proof span 6983-7090 come from the halfA pool card operations/pool-cards/stacks-divisors-halfA.json and were re-read in the pinned chapter file before this card was written.\par\medskip
\noindent\textit{Editorial note.} The proof ends with the author's parenthesis '(some details omitted)' in the glueing paragraph; that phrase is kept verbatim and what it leaves to the reader is the bookkeeping of the closedness argument, while the affine computation and the construction of the module are carried out in the counted span.\par\medskip
\noindent\textit{Editorial note.} Two of the cross-references in this proof point at the lane's own card S13-morph-fpqc-quotient-topology (the quotient topology statement for quasi-compact surjective flat morphisms) and at the openness statement for flat morphisms; both are delivered as navigation pointers.\par\medskip
\noindent\textit{Editorial note.} 11 cross-references occur inside the counted spans; they are listed in external\_references as navigation pointers only.\par\medskip
\noindent\textit{Allowed standard prerequisites (editorial).}\begin{itemize}
\item A flat morphism with integral fibres over an integral base has integral source, cited from the same chapter. Cited at proof lines 6986-6987. Makes the trivializing section of the generic fibre into a regular meromorphic section on the source.
\item Pullbacks of meromorphic functions along a morphism with integral fibres are defined, cited from the same chapter. Cited at proof lines 6988-6990. Allows the product of a pulled back meromorphic function with the chosen section to be considered.
\item Structural facts about a normal local domain of dimension one: it is a discrete valuation ring, the intersection over height one primes with prescribed exponents equals a symbolic power expression, symbolic powers commute with flat base change and flat extensions intersect ideals, and a product of ideals equal to a principal ideal makes the factors finite locally free of rank one. Cited at proof lines 7011-7041. These are the algebra inputs of the affine computation that splits the section s into a product of a unit and an element coming from the base.
\item The quotient topology statement for quasi-compact surjective flat morphisms and the openness statement for flat morphisms locally of finite presentation, cited from the morphisms chapter. Cited at proof lines 7066-7070. Decides the two cases of the hypothesis that the given morphism is either quasi-compact or locally of finite type, showing that the exceptional set is closed.
\end{itemize}\medskip
\ifdefined\hypertarget\hypertarget{sourcepointer-1}{}\fi
\noindent\textit{Source reference.} Original lemma 00PD: focused-seeds/AGP/sources/stacks-snapshot-20261011/tags.txt, source lines 927--927. Same-chapter reference used in the retained author text at source line(s) 1819, 2358, 3021, 7023. It is supplied only as a navigation pointer so that the retained text compiles; the counted item delivers the author's own statement and the author's complete proof, and no part of the referenced item is copied in or substituted for the counted argument.\par\medskip
\ifdefined\hypertarget\hypertarget{sourcepointer-2}{}\fi
\noindent\textit{Source reference.} Original lemma 031S: focused-seeds/AGP/sources/stacks-snapshot-20261011/tags.txt, source lines 3995--3995. Same-chapter reference used in the retained author text at source line(s) 6899, 7019. It is supplied only as a navigation pointer so that the retained text compiles; the counted item delivers the author's own statement and the author's complete proof, and no part of the referenced item is copied in or substituted for the counted argument.\par\medskip
\ifdefined\hypertarget\hypertarget{sourcepointer-3}{}\fi
\noindent\textit{Source reference.} Original lemma 0BBY: focused-seeds/AGP/sources/stacks-snapshot-20261011/tags.txt, source lines 14261--14261. Same-chapter reference used in the retained author text at source line(s) 7046. It is supplied only as a navigation pointer so that the retained text compiles; the counted item delivers the author's own statement and the author's complete proof, and no part of the referenced item is copied in or substituted for the counted argument.\par\medskip
\ifdefined\hypertarget\hypertarget{sourcepointer-4}{}\fi
\noindent\textit{Source reference.} Original lemma 00KV: focused-seeds/AGP/sources/stacks-snapshot-20261011/tags.txt, source lines 765--765. Same-chapter reference used in the retained author text at source line(s) 2994, 3212, 7018. It is supplied only as a navigation pointer so that the retained text compiles; the counted item delivers the author's own statement and the author's complete proof, and no part of the referenced item is copied in or substituted for the counted argument.\par\medskip
\ifdefined\hypertarget\hypertarget{sourcepointer-5}{}\fi
\noindent\textit{Source reference.} Original lemma 031T: focused-seeds/AGP/sources/stacks-snapshot-20261011/tags.txt, source lines 3996--3996. Same-chapter reference used in the retained author text at source line(s) 3417, 6903, 6948, 7032. It is supplied only as a navigation pointer so that the retained text compiles; the counted item delivers the author's own statement and the author's complete proof, and no part of the referenced item is copied in or substituted for the counted argument.\par\medskip
\ifdefined\hypertarget\hypertarget{sourcepointer-6}{}\fi
\noindent\textit{Source reference.} Original lemma 09ME: focused-seeds/AGP/sources/stacks-snapshot-20261011/tags.txt, source lines 12177--12177. Same-chapter reference used in the retained author text at source line(s) 2541, 7053. It is supplied only as a navigation pointer so that the retained text compiles; the counted item delivers the author's own statement and the author's complete proof, and no part of the referenced item is copied in or substituted for the counted argument.\par\medskip
\ifdefined\hypertarget\hypertarget{sourcepointer-7}{}\fi
\noindent\textit{Source reference.} Original lemma 0BC0: focused-seeds/AGP/sources/stacks-snapshot-20261011/tags.txt, source lines 14263--14263. Same-chapter reference used in the retained author text at source line(s) 7045. It is supplied only as a navigation pointer so that the retained text compiles; the counted item delivers the author's own statement and the author's complete proof, and no part of the referenced item is copied in or substituted for the counted argument.\par\medskip
\ifdefined\hypertarget\hypertarget{sourcepointer-8}{}\fi
\noindent\textit{Source reference.} Original lemma 0BCM: focused-seeds/AGP/sources/stacks-snapshot-20261011/tags.txt, source lines 14285--14285. Same-chapter reference used in the retained author text at source line(s) 6989. It is supplied only as a navigation pointer so that the retained text compiles; the counted item delivers the author's own statement and the author's complete proof, and no part of the referenced item is copied in or substituted for the counted argument.\par\medskip
\ifdefined\hypertarget\hypertarget{sourcepointer-9}{}\fi
\noindent\textit{Source reference.} Original lemma 02OU: focused-seeds/AGP/sources/stacks-snapshot-20261011/tags.txt, source lines 3529--3529. Same-chapter reference used in the retained author text at source line(s) 6096, 6992. It is supplied only as a navigation pointer so that the retained text compiles; the counted item delivers the author's own statement and the author's complete proof, and no part of the referenced item is copied in or substituted for the counted argument.\par\medskip
\ifdefined\hypertarget\hypertarget{sourcepointer-10}{}\fi
\noindent\textit{Source reference.} Original lemma 01UA: focused-seeds/AGP/sources/stacks-snapshot-20261011/tags.txt, source lines 2403--2403. Same-chapter reference used in the retained author text at source line(s) 7088. It is supplied only as a navigation pointer so that the retained text compiles; the counted item delivers the author's own statement and the author's complete proof, and no part of the referenced item is copied in or substituted for the counted argument.\par\medskip
\ifdefined\hypertarget\hypertarget{sourcepointer-11}{}\fi
\noindent\textit{Source reference.} Original lemma 02JY: focused-seeds/AGP/sources/stacks-snapshot-20261011/tags.txt, source lines 3353--3353. Same-chapter reference used in the retained author text at source line(s) 7086. It is supplied only as a navigation pointer so that the retained text compiles; the counted item delivers the author's own statement and the author's complete proof, and no part of the referenced item is copied in or substituted for the counted argument.\par\medskip
\begin{lemma}
\label{lemma-in-image-pullback}
Let $\varphi : X \to Y$ be a morphism of schemes. Let $\mathcal{L}$
be an invertible $\mathcal{O}_X$-module. Assume that
\begin{enumerate}
\item $X$ is locally Noetherian,
\item $Y$ is locally Noetherian, integral, and normal,
\item $\varphi$ is flat with integral (hence nonempty) fibres,
\item $\varphi$ is either quasi-compact or locally of finite type,
\item $\mathcal{L}$ is trivial when restricted to the generic fibre of
$\varphi$.
\end{enumerate}
Then $\mathcal{L} \cong \varphi^*\mathcal{N}$ for some invertible
$\mathcal{O}_Y$-module $\mathcal{N}$.
\end{lemma}

\begin{proof}
Let $\xi \in Y$ be the generic point. Let $X_\xi$ be the scheme theoretic
fibre of $\varphi$ over $\xi$. Denote $\mathcal{L}_\xi$ the pullback of
$\mathcal{L}$ to $X_\xi$. Assumption (5) means that $\mathcal{L}_\xi$
is trivial. Choose a trivializing section
$s \in \Gamma(X_\xi, \mathcal{L}_\xi)$. Observe that $X$ is integral by
Lemma \ref{lemma-flat-over-integral-integral-fibre}.
Hence we can think of $s$ as a regular meromorphic section of $\mathcal{L}$.
Pullbacks of meromorphic functions are defined for
$\varphi$ by Lemma \ref{lemma-pullback-meromorphic-sections-defined}.
Let $\mathcal{N} \subset \mathcal{K}_Y$ be the $\mathcal{O}_Y$-module
whose sections over an open $V \subset Y$ are those meromorphic functions
$g \in \mathcal{K}_Y(V)$ such that
$\varphi^*(g)s \in \mathcal{L}(\varphi^{-1}V)$.
A priori $\varphi^*(g)s$ is a section of $\mathcal{K}_X(\mathcal{L})$ over
$\varphi^{-1}V$. We claim that $\mathcal{N}$ is an invertible
$\mathcal{O}_Y$-module and that the map
$$
\varphi^*\mathcal{N} \longrightarrow \mathcal{L},\quad
g \longmapsto gs
$$
is an isomorphism.

\medskip\noindent
We first prove the claim in the following situation:
$X$ and $Y$ are affine and $\mathcal{L}$ trivial. Say $Y = \Spec(R)$,
$X = \Spec(A)$ and $s$ given by the element $s \in A \otimes_R K$
where $K$ is the fraction field of $R$. We can write $s = a/r$
for some nonzero $r \in R$ and $a \in A$. Since $s$ generates $\mathcal{L}$
on the generic fibre we see that there exists an $s' \in A \otimes_R K$
such that $ss' = 1$. Thus we see that $s = r'/a'$ for some
nonzero $r' \in R$ and $a' \in A$. Let
$\mathfrak p_1, \ldots, \mathfrak p_n \subset R$
be the minimal primes over $rr'$. Each $R_{\mathfrak p_i}$
is a discrete valuation ring
(Algebra, Lemmas \ref{algebra-lemma-minimal-over-1} and
\ref{algebra-lemma-criterion-normal}). By assumption
$\mathfrak q_i = \mathfrak p_i A$ is a prime. Hence
$\mathfrak q_i A_{\mathfrak q_i}$ is generated by a single element
and we find that $A_{\mathfrak q_i}$ is a discrete valuation ring as
well (Algebra, Lemma \ref{algebra-lemma-characterize-dvr}). Of course
$R_{\mathfrak p_i} \to A_{\mathfrak q_i}$ has ramification index $1$.
Let $e_i, e'_i \geq 0$ be the  valuation of $a, a'$ in $A_{\mathfrak q_i}$.
Then $e_i + e'_i$ is the valuation of $rr'$ in $R_{\mathfrak p_i}$. Note that
$$
\mathfrak p_1^{(e_1 + e'_1)} \cap \ldots \cap \mathfrak p_i^{(e_n + e'_n)} =
(rr')
$$
in $R$ by Algebra, Lemma
\ref{algebra-lemma-normal-domain-intersection-localizations-height-1}.
Set
$$
I = \mathfrak p_1^{(e_1)} \cap \ldots \cap \mathfrak p_i^{(e_n)}
\quad\text{and}\quad
I' = \mathfrak p_1^{(e'_1)} \cap \ldots \cap \mathfrak p_i^{(e'_n)}
$$
so that $II' \subset (rr')$. Observe that
$$
IA =
(\mathfrak p_1^{(e_1)} \cap \ldots \cap \mathfrak p_i^{(e_n)})A =
(\mathfrak p_1A)^{(e_1)} \cap \ldots \cap (\mathfrak p_i A)^{(e_n)}
$$
by Algebra, Lemmas \ref{algebra-lemma-symbolic-power-flat-extension} and
\ref{algebra-lemma-flat-intersect-ideals}. Similarly for $I'A$. Hence
$a \in IA$ and $a' \in I'A$.
We conclude that $IA \otimes_A I'A \to rr'A$ is surjective.
By faithful flatness of $R \to A$ we find that
$I \otimes_R I' \to (rr')$ is surjective as well.
It follows that $II' = (rr')$ and $I$ and $I'$ are finite locally
free of rank $1$, see
Algebra, Lemma \ref{algebra-lemma-product-ideals-principal}.
Thus Zariski locally on $R$ we can write $I = (g)$ and $I' = (g')$
with $gg' = rr'$. Then $a = ug$ and $a' = u'g'$ for some $u, u' \in A$.
We conclude that $u, u'$ are units. Thus Zariski locally on $R$
we have $s = ug/r$ and the claim follows in this case.

\medskip\noindent
Let $y \in Y$ be a point.
Pick $x \in X$ mapping to $y$. We may apply the result of the previous
paragraph to $\Spec(\mathcal{O}_{X, x}) \to \Spec(\mathcal{O}_{Y, y})$.
We conclude there exists an element $g \in R(Y)^*$ well defined up to
multiplication by an element of $\mathcal{O}_{Y, y}^*$ such that
$\varphi^*(g)s$ generates $\mathcal{L}_x$. Hence $\varphi^*(g)s$
generates $\mathcal{L}$ in a neighbourhood $U$ of $x$.
Suppose $x'$ is a second point lying over $y$ and $g' \in R(Y)^*$ is
such that $\varphi^*(g')s$ generates $\mathcal{L}$ in an open neighbourhood
$U'$ of $x'$. Then we can choose a point
$x''$ in $U \cap U' \cap \varphi^{-1}(\{y\})$
because the fibre is irreducible. By the uniqueness for
the ring map $\mathcal{O}_{Y, y} \to \mathcal{O}_{X, x''}$
we find that $g$ and $g'$ differ (multiplicatively)
by an element in $\mathcal{O}_{Y, y}^*$. Hence we see that $\varphi^*(g)s$
is a generator for $\mathcal{L}$ on an open neighbourhood
of $\varphi^{-1}(y)$. Let $Z \subset X$ be the set of points
$z \in X$ such that $\varphi^*(g)s$ does not generate $\mathcal{L}_z$.
The arguments above show that $Z$ is closed and that $Z = \varphi^{-1}(T)$
for some subset $T \subset Y$ with $y \not \in T$. If we can show that
$T$ is closed, then $g$ will be a generator for $\mathcal{N}$ as an
$\mathcal{O}_Y$-module in the open neighbourhood $Y \setminus T$ of $y$
thereby finishing the proof (some details omitted).

\medskip\noindent
If $\varphi$ is quasi-compact, then $T$ is closed by
Morphisms, Lemma \ref{morphisms-lemma-fpqc-quotient-topology}.
If $\varphi$ is locally of finite type, then $\varphi$ is open
by Morphisms, Lemma \ref{morphisms-lemma-fppf-open}.
Then $Y \setminus T$ is open as the image of the open $X \setminus Z$.
\end{proof}

\end{document}
